A vision-language-action (VLA) model maps perception and instruction
directly to low-level control, in one policy trained end-to-end. The
architecture is elegant, but it inherits the oldest bottleneck of
imitation learning: the policy only knows what its demonstration data
covers. \rev{Moreover,} paired vision-action trajectories are expensive
to collect at the scale a general-purpose manipulator needs. An LLM or
VLM, by contrast, is pretrained on a broader and cheaper corpus: text,
code, and image-text pairs at internet scale. That corpus already
encodes much of what manipulation requires, including object
affordances, spatial and arithmetic reasoning, and precondition
ordering. None of this is robot-specific, and all of it transfers
through code rather than through a demonstration
dataset~\cite{liang2023codeaspolicies, huang2023voxposer,
huang2022innermonologue, singh2023progprompt}. A decoupled
Sense-Plan-Act pipeline makes this architectural bet: the model plans
in language, and a separate execution layer turns the plan into
motion.

A decoupled pipeline still has to ask the model, somewhere, to commit
to numbers: a release height, a grasp width, a target position. Nothing
in an LLM or VLM's pretraining grounds those numbers in the physical
scene the robot acts in. \rev{In simulated validation this cost is
easy to miss: an object's height or the distance a gripper must travel
are never \emph{measured} there, but asserted, correctly, by
construction~\cite{favaliRAS, rohmer2013coppeliasim,
todorov2012mujoco}. In the real world, by contrast, every such
quantity has to be computed: from noisy depth, an imperfect camera, or
a hand-authored scene file that may not match the physical table it
describes.}

This paper closes that gap with an execution harness. We proceeded
in three steps. First, we built an experimental baseline: a ROS2
Sense-Plan-Act stack in which an LLM or VLM decomposes a
natural-language command into a plan of parameterized skills, using any
of six established reasoning strategies \cite{favaliRAS}. The baseline
supports two grounding regimes: an LLM planner with perfect perception,
reading an exact scene description, and a VLM planner perceiving the
scene through a real camera. Second, we ran a structured investigation
of the baseline on pick-and-place tasks with a real UR5cb manipulator.
It surfaced six recurring failure modes, each traced from a hardware
failure back to a specific model output. Crucially, the failures also
occur in the LLM \rev{implementation}, whose perception is perfect by
design. The root cause is therefore how the model reasons about
physical quantities, not how the scene is perceived. Third, guided by
this analysis, we extended the architecture with a third layer between
the high-level planner and the low-level executor: \harness{}
(\harnesslong{}). \rev{The planner keeps the symbolic decisions, which
object to move and where; \harness{} recomputes each action's physical
quantities from a sensor measurement or from an action already
committed in the plan. The result is a complete robotic pipeline that
couples high-level LLM/VLM planning, a deterministic execution
harness, and low-level control primitives. By applying \harness{} we
resolve five of the six failure modes, and a 120-trial real-hardware
benchmark validates the design.}

The contributions of this paper are threefold:
\begin{itemize}
  \item A structured investigation of how LLM/VLM
    planners\rev{, complemented with a set of predefined low-level
    control primitives,} fail on real pick-and-place, yielding a
    taxonomy of six recurring failure modes that arise independently
    of perception quality (Section~\ref{sec:failures}).
  \item \harness{}, a deterministic execution harness for decoupled
    LLM/VLM manipulation planning, presented as an explicit third
    architectural layer between planner and executor. This is the
    paper's main contribution (Section~\ref{sec:harness}).
  \item A 120-trial real-hardware validation of the harness,
    using a single unified model for both grounding modalities. It
    shows that the one failure \harness{} does not yet \rev{resolve},
    multi-object placement collision, dominates the residual failures
    of both \rev{implementations} (Section~\ref{sec:evaluation}).
\end{itemize}

The rest of the paper is organized as follows. The related work
(Section~\ref{sec:related-work}) argues why a decoupled Sense-Plan-Act
architecture with LLM/VLM planners is preferable to an end-to-end VLA,
on data-efficiency and diagnosability grounds, \rev{and reviews the
simulation-based benchmarks this work builds on}. The experimental
baseline (Section~\ref{sec:baseline}) describes the two-layer system
we built, its two perception regimes, and the exact point where the
harness will later be inserted. The failure investigation
(Section~\ref{sec:failures}) reports the six failure modes and the
evidence that they are reasoning failures rather than perception
failures. Section~\ref{sec:harness} introduces \harness{} and shows how
it \rev{resolves} five of the six. The evaluation
(Section~\ref{sec:evaluation}) validates the harness with a 120-trial
benchmark on real hardware. Section~\ref{sec:discussion} discusses
limitations, and Section~\ref{sec:conclusion} concludes.
